% Generated from project outputs. Do not edit by hand.
\begin{table}[!htbp]
\centering
\begin{threeparttable}
\caption{Labor income by subsample: 2021 moments}
\label{tab:annual-2021-y1-moments}
\small
\setlength{\tabcolsep}{4pt}
\renewcommand{\arraystretch}{1.15}
\begin{tabularx}{\linewidth}{@{}Xrrrrr@{}}
\toprule
Sample & N & Mean & SD & Min & Max \\
\midrule
General & 94 & 2,526.1 & 1,509.8 & 286.3 & 6,726.5 \\
Men & 94 & 2,845.5 & 1,577.7 & 398.3 & 8,403.7 \\
Women & 94 & 2,023.7 & 1,576.3 & 47.1 & 7,548.6 \\
Urban & 94 & 2,810.3 & 1,556.3 & 254.1 & 6,445.8 \\
Rural & 94 & 2,088.0 & 1,303.4 & 250.0 & 7,380.5 \\
Indigenous & 94 & 2,168.6 & 1,742.8 & 238.3 & 9,778.0 \\
Non-Indigenous & 94 & 2,701.2 & 1,480.6 & 104.9 & 6,945.8 \\
Bottom 99 percent & 94 & 2,347.9 & 1,310.1 & 286.3 & 6,214.5 \\
Bottom 90 percent & 94 & 1,940.7 & 871.7 & 309.1 & 3,838.8 \\
Bottom 50 percent & 94 & 1,073.0 & 412.2 & 153.6 & 2,307.3 \\
\bottomrule
\end{tabularx}
\begin{tablenotes}[flushleft]
\footnotesize
\item Monthly quetzales. Unweighted distribution of survey-weighted cohort means before transition matching. N counts cohort cells, not individuals. SD is dispersion, not a standard error. Subsamples overlap. Cumulative income definitions and treatment of missing components follow the merging scripts.
\end{tablenotes}
\end{threeparttable}
\end{table}
